% ID: SEL-D0011-BC-E06
% ESTADO: validada. S012; tres pasadas del mismo autor, sin segundo revisor.
\textbf{a) Complemento e intersección.}\par\nopagebreak[4]
\[
\boxed{P(A)=1-0.4=0.6.}
\]
Como $A$ es unión disjunta de $A\cap B$ y $A\cap B^c$,
\[
\boxed{P(A\cap B)=0.6-0.12=0.48.}
\]
\textbf{b) Condición de independencia.}\par\nopagebreak[4]
La independencia equivale a $P(A\cap B)=P(A)P(B)$. Puesto que $P(A)>0$,
\[
0.48=0.6P(B)\quad\Longrightarrow\quad\boxed{P(B)=0.8.}
\]
Es un valor posible: las cuatro regiones disjuntas tendrían probabilidades $0.48,0.12,0.32,0.08$, todas no negativas y de suma uno.

\textbf{c) Tres probabilidades con el dato adicional.}\par\nopagebreak[4]
Ahora se da $P(B^c)=0.2$, luego $P(B)=0.8$. No es necesario trasladar una hipótesis de independencia desde b: los datos y las identidades bastan.
Por inclusión-exclusión,
\[
\boxed{P(A\cup B)=0.6+0.8-0.48=0.92.}
\]
Por De Morgan,
\[
\boxed{P(A^c\cup B^c)=1-P(A\cap B)=0.52.}
\]
La notación $P(A/B^c)$ del original significa condicionada:
\[
\boxed{P(A\mid B^c)=\frac{P(A\cap B^c)}{P(B^c)}=\frac{0.12}{0.2}=0.6.}
\]
Se distingue la unión de los complementarios del complemento de la unión, que valdría $0.08$.
